\documentclass[10pt,a4paper]{amsart}
\usepackage[margin=19mm]{geometry}
\usepackage{amsmath,amssymb,mathtools}
\usepackage[scr=rsfs,cal=euler]{mathalfa}
\usepackage{xfrac}
\usepackage[unicode,hidelinks,bookmarks=false]{hyperref}
\newcommand{\ie}{i.e.\ }
\newcommand{\cf}{cf.\ }
\newcommand{\resp}{respectively\ }
\newcommand{\Subsection}{\subsection}
\providecommand{\nobreakdash}{\nobreak\hbox{-}}
\setlength{\emergencystretch}{2em}
\multlinegap0pt
\usepackage{mathtools}

\theoremstyle{plain}
\newtheorem{Thm}{Theorem}[section]
\newtheorem{Que}[Thm]{Question}
\newtheorem{Lem}[Thm]{Lemma}
\newtheorem{Prop}[Thm]{Proposition}
\newtheorem{Conj}[Thm]{Conjecture}
\newtheorem{Cor}[Thm]{Corollary}
\newtheorem{Cla}[Thm]{Claim}
\newtheorem{Fac}[Thm]{Fact}
\theoremstyle{definition}
\newtheorem{Axi}[Thm]{Axiom}
\newtheorem{Ass}[Thm]{Assumption}
\newtheorem{Def}[Thm]{Definition}
\newtheorem{Rem}[Thm]{Remark}
\newtheorem{Exa}[Thm]{Example}
\title[Cutoff on non-negatively curved chains without symmetric support condition]{Cutoff on non-negatively curved chains without symmetric support condition}
\author{Kazuki Okamura}
\date{Source: arXiv:2609.28518v1 (CC BY 4.0)}
\begin{document}
\maketitle

\noindent\emph{Source and licence.} The delivered work is the article shown in the title above, version arXiv:2609.28518v1 (CC BY 4.0), distributed under the Creative Commons Attribution 4.0 International licence, as recorded in the arXiv licence metadata for this exact version and shown on the abstract page saved with this item. The full licence notice, which carries the licence URL and the address of that abstract page, is the bundled file \texttt{licenses/arxiv-2609.28518-CC-BY-4.0.txt}.
\par\smallskip

\noindent\textit{Editorial scope.} Counted unit: the article's Prop (source label prop:6.3) (source0.tex lines 489--492, source label \texttt{prop:6.3}): ``For every , diam ( X ) (1 k ( X , P) ) (t mix ( ) 2t mix ( ) 1- k ( X , P) t rel 1- ).''. It is delivered together with the article's complete original proof (source0.tex lines 495--564, one proof environment printed later in the article, detached from the statement by the article's own arrangement, 70 lines), copied line for line from the article's own text. Required context retained uncounted: eq:def-relax (lines 158--160). Editorial formatting only: an AMS \texttt{amsart} preamble that restates the article's own macro definitions and its own theorem declarations, and the theorem environments the retained lines use; the packages mathtools are loaded to match the article's own setup; the article's own class file is neither shipped nor input; the retained environments are renumbered sequentially in order of appearance, whereas the article prints them with its own numbering; the article's equation numbering is not reproduced and every cross-reference inside the retained text resolves to the wrapper's numbering. No citation occurs inside any retained line, so the wrapper carries no bibliography; All retained bodies are exact copies of the bundled \texttt{sources/211597/source0.tex}, so no omission receipt applies. No asset-inclusion command occurs inside any retained span and the package ships no image asset.


\section*{Retained context: eq:def-relax (source0.tex lines 158--160)}

\begin{equation}\label{eq:def-relax} 
t_{\mathrm{rel}} \coloneqq  \sup \left\{ \frac{\mathrm{Var}_\pi(f)}{\mathcal{E}(f)} \middle| f : \mathcal{X} \to \mathbb{R},\ \mathcal{E}(f) > 0 \right\}. 
\end{equation}


\section{The counted unit: Prop (source label prop:6.3) and its complete original proof}

\begin{Prop}\label{prop:6.3}
For every $\varepsilon \in (0,1)$,
\[ \mathrm{diam}(\mathcal{X}) \le (1+k_{(\mathcal{X}, P)})\left(t_{\mathrm{mix}}(\varepsilon)+\sqrt{\frac{2t_{\mathrm{mix}}(\varepsilon)}{1-\varepsilon}}+k_{(\mathcal{X}, P)}\sqrt{\frac{t_{\mathrm{rel}}}{1-\varepsilon}}\right).\]
\end{Prop}

\begin{proof}
Fix $o \in \mathcal{X}$. 
Let $t \coloneqq t_{\mathrm{mix}}(\varepsilon) \ge 0$. 
Let
\[ A_{o,t,\varepsilon} \coloneqq \left\{x\in\mathcal{X}\ \middle|\ d_P (o \to x) \le t+\sqrt{\frac{2t}{1-\varepsilon}}\right\}. \]

Let $(Y_n)_{n \ge 0}$ be a Markov chain on $\mathcal{X}$ with transition probability $P$ and $Y_0 = o$,
and let $(N_t)_{t \ge 0}$ be a Poisson process with rate $1$ which is independent of $(Y_n)_{n \ge 0}$.
Then the distribution of $Y_{N_t}$ is $\mathscr{P}_t (o, \cdot)$.
Since $P(Y_{n}, Y_{n+1}) > 0$ almost surely for every $n \ge 0$, we see that $d_P (o \to Y_n) \le n$ almost surely for every $n \ge 0$,
and hence $d_P (o \to Y_{N_t}) \le N_t$ almost surely;
that is, the distribution of $d_P (o \to Y_{N_t})$ is stochastically dominated by the Poisson distribution with mean $t$.

Hence
\[ \mathscr{P}_t\left(o, A_{o,t,\varepsilon}^c\right) = \mathbb{P} \left(Y_{N_t} \notin A_{o,t,\varepsilon}\right) \]
\[ = \mathbb{P} \!\left(d_P (o \to Y_{N_t}) > t+\sqrt{\frac{2t}{1-\varepsilon}}\right) \le \mathbb{P}\!\left(N_t >  t+\sqrt{\frac{2t}{1-\varepsilon}}\right). \]

Assume that $t > 0$. 
Since $\mathbb{E}[N_t] = \mathrm{Var}(N_t) = t > 0$, by the Chebyshev inequality,
\[ \mathbb{P}\!\left(N_t > t+\sqrt{\frac{2t}{1-\varepsilon}}\right) \le  \frac{1-\varepsilon}{2}. \]
Therefore, 
\begin{equation}\label{eq:Pt-lower-annulus}
\mathscr{P}_t \left(o,A_{o,t,\varepsilon}\right)\ge \frac{1+\varepsilon}{2}. 
\end{equation}
If $t=0$, then $A_{o,0,\varepsilon} = \{o\}$ and $\mathcal P_0(o,A_{o,0,\varepsilon})=1$, and hence, \eqref{eq:Pt-lower-annulus} holds. 


Since $t = t_{\mathrm{mix}}(\varepsilon)$, we obtain that
$\mathscr{P}_t \left(o,A_{o,t,\varepsilon}\right)\le \pi(A_{o,t,\varepsilon})+\varepsilon$. 
By this and \eqref{eq:Pt-lower-annulus}, it holds that 
\[ \pi\left(A_{o,t,\varepsilon}\right) \ge \frac{1-\varepsilon}{2}.\]

Let $f(x)\coloneqq d_P (o \to x)$ for $x \in \mathcal{X}$, and let $U$ be a random variable with values in $\mathcal{X}$ whose distribution is $\pi$.
Then
\begin{equation}\label{eq:dist-half}
\mathbb{P} \left(f(U)\le t+\sqrt{\frac{2t}{1-\varepsilon}}\right) = \pi\left(A_{o,t,\varepsilon}\right) \ge \frac{1-\varepsilon}{2}.
\end{equation}

We remark that for every $x \in \mathcal{X}$, $2\Gamma(f,f)(x) \le \|f\|_{\mathrm{Lip}}^2$.
Hence, $2\mathcal{E}(f) \le \|f\|_{\mathrm{Lip}}^2$.
By \eqref{eq:def-relax} and the estimate that $\|f\|_{\mathrm{Lip}} \le k_{(\mathcal{X}, P)}$,
we obtain that
\[ \mathrm{Var}_\pi (f)\le t_{\mathrm{rel}}\mathcal{E}(f) \le \frac{1}{2}\, t_{\mathrm{rel}} \left\|f \right\|_{\mathrm{Lip}}^2 \le \frac{1}{2}\, t_{\mathrm{rel}} k_{(\mathcal{X}, P)}^2. \]

By this estimate and Chebyshev's inequality, we see that for every $\varepsilon^{\prime} \in (\varepsilon, 1)$, 
\begin{equation*}
\mathbb{P}\left(\left|f(U)- \mathbb{E}[f(U)]\right|\le k_{(\mathcal{X}, P)}\sqrt{\frac{t_{\mathrm{rel}}}{1-\varepsilon^{\prime} }} \right)\ge \frac{1+\varepsilon^{\prime} }{2}.
\end{equation*}

By this and \eqref{eq:dist-half}, 
\[ \left\{f(U)\le t_{\mathrm{mix}}(\varepsilon) + \sqrt{\frac{2 t_{\mathrm{mix}}(\varepsilon)}{1-\varepsilon}}\right\} \cap \left\{\left|f(U)- \mathbb{E}[f(U)]\right|\le k_{(\mathcal{X}, P)}\sqrt{\frac{t_{\mathrm{rel}}}{1-\varepsilon^{\prime} }}\right\} \ne \emptyset. \]

Therefore,
\[  \mathbb{E}\left[d_P (o \to U) \right] = \mathbb{E}[f(U)] \le t_{\mathrm{mix}}(\varepsilon) + \sqrt{\frac{2t_{\mathrm{mix}}(\varepsilon)}{1-\varepsilon}} +k_{(\mathcal{X}, P)}\sqrt{\frac{t_{\mathrm{rel}}}{1-\varepsilon^{\prime} }}. \]
By letting $\varepsilon^{\prime} \to \varepsilon$,
we obtain that
\begin{equation}\label{eq:exp-dist-pi}
\mathbb{E}\left[d_P (o \to U) \right] \le t_{\mathrm{mix}}(\varepsilon) + \sqrt{\frac{2t_{\mathrm{mix}}(\varepsilon) }{1-\varepsilon}} +k_{(\mathcal{X}, P)}\sqrt{\frac{t_{\mathrm{rel}}}{1-\varepsilon}}.
\end{equation}

We see that
$d_P (x \to y) \le k_{(\mathcal{X}, P)}\,d_P (y \to x)$ for $x, y \in \mathcal{X}$,
by induction on the value of $d_P (y\to x)$.
Hence, for every $o^{\prime}, x \in \mathcal{X}$,
\[ d_P (o \to o^{\prime}) \le d_P (o\to x) + d_P (x\to o') \le d_P (o\to x)+k_{(\mathcal{X}, P)}\,d_P (o^{\prime} \to x). \]

By this and \eqref{eq:exp-dist-pi},
\[ d_P (o \to o') \le \mathbb{E}[d_P (o\to U)] + k_{(\mathcal{X}, P)}\, \mathbb{E}\left[d_P (o^{\prime}  \to U)\right] \]
\[ \le (k_{(\mathcal{X}, P)}+1)\left(t_{\mathrm{mix}}(\varepsilon) + \sqrt{\frac{2 t_{\mathrm{mix}}(\varepsilon)}{1-\varepsilon}} +k_{(\mathcal{X}, P)}\sqrt{\frac{t_{\mathrm{rel}}}{1-\varepsilon}} \right). \qedhere \]
\end{proof}

\end{document}
